\section{Ferramentas para olhar a JVM}
\label{sec:ferramentas}

O \texttt{javap} mostra o que o \texttt{javac} gerou. Para ver o que a
JVM faz com isso em execução, estas são as ferramentas, todas com
documentação oficial:

\small
\begin{longtable}{@{}p{3.2cm} p{9cm}@{}}
\toprule
Ferramenta & Para que serve \\
\midrule \endhead
\texttt{javap} & Desmonta \texttt{.class}: \texttt{-c}, \texttt{-p},
  \texttt{-v} (\citetitle{doc:javap}) \\
\texttt{-XX:+\allowbreak PrintCompilation} & Uma linha a cada método compilado pelo
  JIT (\citetitle{doc:java-cmd}) \\
\texttt{-XX:+\allowbreak LogCompilation} & Log detalhado de compilação em
  \texttt{hotspot\_pid<pid>.log} (o manual ainda diz \texttt{hotspot.log});
  exige \texttt{-XX:+UnlockDiagnosticVMOptions} \\
JITWatch & Analisador e visualizador do log de compilação da HotSpot:
  inlining, métodos quentes, bytecode e assembly
  (\citetitle{doc:jitwatch}) \\
JFR & Flight Recorder: grava eventos da aplicação, da JVM e do sistema
  operacional num arquivo, com meta de no máximo 1\% de overhead
  (\citetitle{doc:jep328}, JDK~11). Inicie com
  \texttt{-XX:StartFlightRecording} ou \texttt{jcmd <pid> JFR.start} \\
\texttt{jfr} & Lê gravações: \texttt{print}, \texttt{summary},
  \texttt{view} (\citetitle{doc:jfr-tool}) \\
\texttt{jcmd} & Comandos para uma JVM rodando: \texttt{JFR.start},
  \texttt{Thread.dump\_to\_file}, \texttt{GC.heap\_info},
  \texttt{VM.flags}, \texttt{Compiler.queue} (\citetitle{doc:jcmd}) \\
JMC & JDK Mission Control, projeto do OpenJDK com ferramentas gráficas
  que abrem gravações do JFR (\citetitle{doc:jmc}) \\
async-profiler & Profiler por amostragem de CPU e alocação, que segundo
  a documentação não sofre de \emph{safepoint bias}; gera flame graphs
  (\citetitle{doc:async-profiler}) \\
JMH & Harness para microbenchmarks na JVM (\citetitle{doc:jmh}) \\
\bottomrule
\end{longtable}
\normalsize

\paragraph{Por onde começar.}
Para \enquote{está lento em produção}, comece pelo JFR: o overhead é
baixo, ele já vem no JDK e cobre GC, alocação, locks e CPU. Para
\enquote{este trecho ficou mais rápido?}, use JMH, nunca um laço com
\texttt{nanoTime}. Para \enquote{por que o JIT não otimizou este
método?}, \texttt{PrintCompilation} e depois JITWatch. Para \enquote{o
que exatamente o \texttt{javac} gerou?}, \texttt{javap}.
